% Anonymizace a pseudonymizace: metody a limity
\textbf{General background knowledge:} Anonymizace = odstranění nebo transformace údajů tak, aby subjekt nebyl identifikovatelný; pseudonymizace = nahrazení identifikátorů stabilním pseudonymem, obnova identity možná jen s klíčem.
Praktická doporučení: princip minimalizace — sbírat jen nezbytná pole; místo jmen/fotek používat pseudonymy/ID \cite{kb_ai_101_tipu_pdf_04e4a115_page_8_1}.
Před sdílením s externími nástroji aplikujte anonymizaci/pseudonymizaci a omezte přístupy; neuploadujte studentské fotografie do neověřených cloudů \cite{kb_ai_101_tipu_pdf_04e4a115_page_8_1,kb_ai_101_tipu_pdf_04e4a115_page_53_2}.
Rychlý pracovní postup:
\begin{enumerate}
\item Definujte účel a uchovávejte jen nezbytná pole; pokud není třeba identifikovat, neukládejte ID \cite{kb_ai_101_tipu_pdf_04e4a115_page_8_1}.
\item Před sdílením zkontrolujte PII, aplikujte anonymizaci/pseudonymizaci a odstraňte citlivá metadata \cite{kb_ai_101_tipu_pdf_04e4a115_page_8_1}.
\item Ukládejte šifrovaně v interním úložišti, nastavte přístupy; v pilotech ověřte DPIA/DPA a retenci \cite{kb_ai_101_tipu_pdf_04e4a115_page_8_1,kb_malinka_chatgpt_ve_vyuce_dva_roky_pote_2024_pdf_90bc3aaf_page_12_1}.
\end{enumerate}
Krátké příklady: kvíz — anonymní skóre nebo interní ID; programátorský projekt — školní repozitář a minimalizace PII v commitech; sdílení s externí AI — sjednat odstranění PII a rozsah v DPA \cite{kb_ai_101_tipu_pdf_04e4a115_page_8_1,kb_ai_101_tipu_pdf_04e4a115_page_53_2,kb_malinka_chatgpt_ve_vyuce_dva_roky_pote_2024_pdf_90bc3aaf_page_12_1}.
Limity a checklist: anonymizace není absolutní — posuzujte reziduální riziko a DPIA; držte mapovací tabulky pseudonymů odděleně; před sdílením ověřte nezbytnost polí, pseudonymizaci, odstranění metadat/fotek a pro piloty DPIA/DPA s lidským dohledem.